\section{ĐỀ SỐ \thesection }%ĐỀ 6
\setcounter{bt}{0}

%%Câu 1.
\begin{bt} \textbf{ } 
	\begin{listEX}
		\item Em hãy tìm: ƯCLN$(108,60)$.
		\item Sắp xếp các số nguyên sau theo thứ tự giảm dần: $-20\,623 ;\ 0 ;\ 15 ;\ -20 ;\ 7 ;\ -99$
		\item Một kho đông lạnh đang ở nhiệt độ là $-5^{\circ}$C. Sau đó, một công nhân đã bấm nút điều chỉnh tăng nhiệt độ trong kho lên $8^{\circ} \mathrm{C}$. Hỏi nhiệt độ trong kho đông lạnh lúc này là bao nhiêu?
	\end{listEX}
	
	\loigiai{a) ƯCLN$(108, 60) = 12$ \hfill b) $15, 7, 0, -20, -99, -2023$ \hfill c) $3^{\circ} C$
	}
\end{bt}


%%Câu 2. 
\begin{bt} \textbf{ }
	\begin{listEX}[2]
		\item $43+(-100)+(-43)+150$;
		\item $5^{2022}:5^{2019}-[130-(12-2.3)^2]$;
		\item $S = -1+3-5+7-9+11-...-201+203$.
	\end{listEX}
	\loigiai{a) $50$ \hfill b) $31$ \hfill c) $202$.
	}
\end{bt}




%%Câu 3.
\begin{bt} Tìm số tự nhiên $x$ biết:
	\begin{listEX}[3]
		\item $2x + 15 = 57$;
		\item $100 - (x + 15) = 150$;
		\item $x \in \text{B}(12)$ và $40 < x \leqslant 50$.
	\end{listEX}
	
	\loigiai{a) $21$ \hfill b) $-65$ \hfill c) $48$}
\end{bt}


%%Câu 4.
\begin{bt} Một trường trung học cơ sở trên địa bàn TP HCM tổ chức cho tất cả các em học sinh của mình đi tham quan, hướng nghiệp. Khi xếp hàng $8$ em, hàng $10$ em, hàng $12$ em để học sinh tham gia các trò chơi trải nghiệm ngoài trời thì không thừa và không thiếu bạn nào. Hỏi trường đó có tất cả bao nhiêu học sinh, biết rằng tổng số học sinh của trường đó khoảng từ $400$ đến $500$ em.
	
	\loigiai{$480$ học sinh}
\end{bt}


%%Câu 5.
\begin{bt} Một cửa hàng A thống kê số lượng cam bán được trong $4$ ngày cuối tuần vừa rồi bằng biểu đồ tranh sau đây:
	\begin{center}
		\def\icon{\tikz{
				\fill[orange] (0,0) circle (0.4);
				\fill[white] (0,0) circle (0.35) (0,0) circle (0.2cm);
				\fill[orange!60!yellow] (0,0) circle (0.29);
				\foreach \angle in {0,45,...,315} { \draw[white, line width=1pt] (0,0) -- (\angle:0.29);}
		}}
		\def\hicon{\tikz{\clip(-0.4,-0.4)rectangle(0,0.4);\fill[orange] (0,0) circle (0.4);
				\fill[white] (0,0) circle (0.35) (0,0) circle (0.2cm);
				\fill[orange!60!yellow] (0,0) circle (0.29);
				\foreach \angle in {0,45,...,315} { \draw[white, line width=1pt] (0,0) -- (\angle:0.29);}}}
		\def\vetranh#1{\foreach \dem in{1,...,#1}{\icon \ }}{
			\begin{tabular}{|c|p{7cm}<{}|}\hline 
				\multicolumn{2}{|c|}{\bf Số lượng cam bán được trong 4 ngày cuối tuần vừa rồi của cửa hàng A }\\  \hline 
				Thứ Năm &\vetranh{5}\\ \hline 
				Thứ Sáu&\vetranh{3}\ \hicon  \\ \hline 
				Thứ Bảy&\vetranh{4} \\ \hline 
				Chủ nhật &\vetranh{7} \\ \hline 
			\end{tabular}\\
			Biết rằng: \icon = 10 kg;        \hicon = 5 kg}
	\end{center}     
	\begin{enumerate}
		\item Lập bảng thống kê tương ứng với biểu đồ tranh trên.
		\item Từ bảng thống kê, em hãy cho biết: 
		\begin{itemize}
			\item Thứ mấy cửa hàng đó bán được ít cam nhất?
			\item Tổng số cam bán được trong 4 ngày cuối tuần vừa rồi là bao nhiêu kg?
		\end{itemize} 
	\end{enumerate}
	
	
	
	
	\loigiai{a) Học sinh tự lập bảng \hfill b) Thứ $6$; $195$kg
	}
\end{bt}

%%Câu 6.
\begin{bt} 	Một khu đất của ông Tám có dạng như hình vẽ bên:
	\immini{			
		\begin{listEX}
			\item Tính chu vi của khu đất nhà ông Tám.
			\item Ông Tám mua thêm một miếng đất hình vuông liền kề với khu đất của ông có kích thước cạnh bằng $400$ cm. Sau đó, ông dùng những viên gạch hình bình hành lát nền toàn bộ đất của ông với chi phí mua gạch là $7$ triệu đồng. Em hãy tính diện tích một viên gạch hình bình hành này là bao nhiêu dm$^{2}$, biết rằng giá $1$ thùng gạch (gồm $12$ viên) là $168\,000$ đồng và khoảng cách giữa các viên gạch là không đáng kể.
		\end{listEX}
	}{
		\begin{tikzpicture}[scale=1,line width=0.8pt,line join=round,>=stealth]
			\def\a{2}
			\path (0,0) coordinate (F)
			(0,6/\a) coordinate (A)
			(7/\a,6/\a) coordinate (B)
			(7/\a,4/\a) coordinate (C)
			(10/\a,4/\a) coordinate (D)
			(10/\a,0) coordinate (E)
			;
			\draw (A)--(B)node[midway,above]{7 m}
			--(C)node[midway,left]{2 m}
			--(D)node[midway,below]{3 m}
			--(E)node[midway,right]{4 m}
			--(F)node[midway,below]{10 m}
			--(A)node[midway,left]{6 m}
			;
			\foreach \x/ \goc in {A/135,B/45,C/45,D/45,E/-45,F/-135}  \fill (\x) circle (0pt)        
			($(\x)+(\goc:3.5mm)$) node {$\x$};
		\end{tikzpicture}
	}
	\loigiai{a) $32$ m \hfill b) $14$ dm$^2$
	}
\end{bt}


%%Câu 7.
\begin{bt} Tìm số nguyên tố $n$ sao cho $(3n+7)$ là bội của $(n+1)$.   
	
	\loigiai{$n = 3$}
\end{bt}